\documentclass[11pt,a4paper]{article}
\usepackage[utf8]{inputenc}
\usepackage[T1]{fontenc}
\usepackage[catalan]{babel}
\usepackage{lmodern}
\usepackage[sfdefault,lining]{FiraSans}
\usepackage[a4paper,top=2.3cm,bottom=2.2cm,left=1.8cm,right=1.8cm,headheight=1.3cm,headsep=0.4cm,footskip=1cm]{geometry}
\usepackage{amsmath,amssymb}
\usepackage{icomma}
\usepackage{graphicx}
\usepackage[table]{xcolor}
\usepackage{tikz}
\usepackage{fontawesome5}
\usepackage{ccicons}
\usepackage{fancyhdr}
\usepackage{lastpage}
\usepackage{tasks}
\usepackage{needspace}
\usepackage[most]{tcolorbox}
\usepackage{array}
\usepackage[hidelinks,pdftitle={Solucions de la fitxa de la unitat 4: Límits i continuïtat. 2n de Batxillerat CCSS},pdfsubject={Matemàtiques Aplicades a les Ciències Socials II},pdfkeywords={límits, límits a l'infinit, comparació d'infinits, indeterminacions, límits laterals, continuïtat, funcions a trossos, paràmetres, asímptotes, PAU}]{hyperref}

\definecolor{navy}{HTML}{1B2748}
\definecolor{accent}{HTML}{E8603A}
\definecolor{teal}{HTML}{1F9E90}
\definecolor{soft}{HTML}{EEF1F8}
\definecolor{muted}{HTML}{5D6685}

\setlength{\parindent}{0pt}
\setlength{\parskip}{3pt}
\settasks{label=\alph*), label-format=\bfseries\color{navy}, label-width=1.3em, item-indent=1.8em, column-sep=0.6em, after-item-skip=5pt, before-skip=3pt, after-skip=4pt}

% ---------- capçalera i peu ----------
\newcommand{\llicencia}{\href{https://creativecommons.org/licenses/by-nc-sa/4.0/deed.ca}{\ccbyncsa}\ \ Llicència Creative Commons Reconeixement-NoComercial-CompartirIgual 4.0 Internacional (CC BY-NC-SA 4.0)}
\pagestyle{fancy}
\fancyhf{}
\fancyhead[L]{\raisebox{-0.35\height}{\includegraphics[height=0.95cm]{logo-gva.png}}}
\fancyhead[C]{\small\color{muted}Solucions de la fitxa · Unitat 4: Límits i continuïtat · 2n Batxillerat CCSS}
\fancyhead[R]{\raisebox{-0.35\height}{\includegraphics[height=0.95cm]{logo-vermellar.png}}}
\fancyfoot[L]{\scriptsize\color{muted}\llicencia}
\fancyfoot[R]{\small\color{muted}\thepage}
\renewcommand{\headrulewidth}{0.4pt}
\renewcommand{\headrule}{\color{muted!40}\hrule height \headrulewidth}
\fancypagestyle{primera}{\fancyhf{}\renewcommand{\headrulewidth}{0pt}
  \fancyfoot[L]{\scriptsize\color{muted}\llicencia}
  \fancyfoot[R]{\small\color{muted}\thepage}}

% ---------- nivells i exercicis ----------
\newcommand{\estrelles}[1]{{\color{accent}\ifcase#1\or\faStar\faStar[regular]\faStar[regular]\or\faStar\faStar\faStar[regular]\or\faStar\faStar\faStar\fi}}
\newcounter{exer}
\newcommand{\ex}[1]{\par\needspace{4\baselineskip}\medskip\refstepcounter{exer}\noindent{\color{navy}\bfseries\theexer.}\ \ }

\newcommand{\seccio}[2]{\par\needspace{10\baselineskip}\bigskip
  \begin{tcolorbox}[colback=navy,colframe=navy,coltext=white,arc=2mm,boxrule=0pt,left=3mm,right=3mm,top=1.2mm,bottom=1.2mm]
  {\large\bfseries #1\quad #2}
  \end{tcolorbox}\vspace{-2pt}}


% quadre de recordatori (només a la fitxa d'exercicis)
\newcommand{\recorda}[1]{\par\smallskip\begin{tcolorbox}[colback=accent!7,colframe=accent!55,arc=2mm,boxrule=0.5pt,left=2.5mm,right=2.5mm,top=1mm,bottom=1mm]\small{\color{accent}\faLightbulb[regular]}\ \textbf{Recorda:}\ #1\end{tcolorbox}}
% etiqueta dels exercicis de la PAU d'altres comunitats
\newcommand{\pau}[1]{{\color{teal}\small\bfseries\faUniversity\ PAU · #1.}\ }
\newcommand{\res}[1]{\boxed{#1}}\newcommand{\rest}[1]{\fbox{#1}}

% =====================================================================
% Capçalera FentMates de les fitxes: un color per a cada curs.
% Els colors són els de la pàgina inicial (LEVELS a assets/js/materials.js);
% el logo ve de _eines/logo/pdf/ (python3 _eines/logo/genera.py).
%
% Ús, al final del preàmbul de la unitat:
%   % =====================================================================
% Capçalera FentMates de les fitxes: un color per a cada curs.
% Els colors són els de la pàgina inicial (LEVELS a assets/js/materials.js);
% el logo ve de _eines/logo/pdf/ (python3 _eines/logo/genera.py).
%
% Ús, al final del preàmbul de la unitat:
%   % =====================================================================
% Capçalera FentMates de les fitxes: un color per a cada curs.
% Els colors són els de la pàgina inicial (LEVELS a assets/js/materials.js);
% el logo ve de _eines/logo/pdf/ (python3 _eines/logo/genera.py).
%
% Ús, al final del preàmbul de la unitat:
%   \input{../comu/fentmates.tex}
%   \fmcurs{1eso}
% i, al començament del document:
%   \capcalera{Matemàtiques}{1r ESO}{Fitxa d'exercicis · Unitat 2}{Divisibilitat}
% Opcional abans de \input: \def\fmneed{14} (espai mínim abans de cada secció).
% =====================================================================
\makeatletter
\def\fm@logodir{../../logo/pdf}
\providecommand{\fmneed}{10}
\def\fm@cursnom{}

% clau del curs, color i color del text damunt del color
\newcommand{\fm@curs}[3]{\@namedef{fm@color@#1}{#2}\@namedef{fm@text@#1}{#3}}
\fm@curs{1eso}{E6B84A}{navy}
\fm@curs{2eso}{3B82F6}{white}
\fm@curs{3eso}{FF7A50}{white}
\fm@curs{4eso-a}{A78BFA}{white}
\fm@curs{4eso-b}{6D28D9}{white}
\fm@curs{1batx-mi}{9B1C1C}{white}
\fm@curs{1batx-ccss}{EF6B6B}{white}
\fm@curs{2batx-mii}{1E6B47}{white}
\fm@curs{2batx-ccss}{2BB3A3}{white}

\newcommand{\fmcurs}[1]{%
  \@ifundefined{fm@color@#1}{\PackageError{fentmates}{Curs desconegut: #1}{}}{}%
  \edef\fm@tmp{\noexpand\definecolor{curs}{HTML}{\@nameuse{fm@color@#1}}%
               \noexpand\colorlet{curstext}{\@nameuse{fm@text@#1}}}\fm@tmp
  \colorlet{cursfosc}{curs!72!black}%
  \def\fm@logo{\fm@logodir/vertical-#1.pdf}%
  \def\fm@marca{\fm@logodir/marca-#1.pdf}}

% ---------- capçalera de la primera pàgina ----------
% #1 matèria (petit), #2 curs, #3 tipus de document i unitat, #4 títol
\newcommand{\capcalera}[4]{%
  \gdef\fm@cursnom{#2}%
  \par\noindent
  \begin{tcolorbox}[colback=curs,colframe=curs,boxrule=0pt,arc=2.2mm,width=3.5cm,height=2.15cm,
      valign=center,left=3mm,right=2mm,top=1mm,bottom=1mm,nobeforeafter,box align=center]
    \color{curstext}{\footnotesize #1}\par\vspace{1pt}{\fontsize{19}{21}\selectfont\bfseries #2\par}%
  \end{tcolorbox}\hspace{5mm}%
  \begin{minipage}[c]{\dimexpr\linewidth-3.5cm-5mm-2.9cm\relax}
    {\color{cursfosc}\small\bfseries\MakeUppercase{#3}}\par\vspace{2pt}
    {\color{navy}\fontsize{20}{23}\selectfont\bfseries #4\par}
  \end{minipage}\hfill
  \begin{minipage}[c]{2.6cm}\raggedleft\includegraphics[width=2.6cm]{\fm@logo}\end{minipage}\par
  \vspace{2pt}}

% ---------- barres de secció del color del curs ----------
\renewcommand{\seccio}[2]{\par\needspace{\fmneed\baselineskip}\bigskip
  \begin{tcolorbox}[colback=curs,colframe=curs,coltext=curstext,arc=2mm,boxrule=0pt,left=3mm,right=3mm,top=1.2mm,bottom=1.2mm]
  {\large\bfseries #1\quad #2}
  \end{tcolorbox}\vspace{-2pt}}

% ---------- capçalera i peu de les pàgines ----------
% primera pàgina: logos institucionals al peu
\fancypagestyle{primera}{\fancyhf{}\renewcommand{\headrulewidth}{0pt}
  \fancyfoot[L]{\raisebox{-0.32\height}{\includegraphics[height=0.78cm]{logo-gva.png}}\hspace{5mm}%
    \raisebox{-0.32\height}{\includegraphics[height=0.78cm]{logo-vermellar.png}}}
  \fancyfoot[R]{\scriptsize\color{muted}\href{https://creativecommons.org/licenses/by-nc-sa/4.0/deed.ca}{\ccbyncsa}\ \ CC BY-NC-SA 4.0 · fentmates.org\qquad\small\thepage}}
% la resta: el símbol a l'esquerra i el curs a la dreta, amb el títol al centre com abans
\fancyhead[L]{\raisebox{-0.3\height}{\includegraphics[height=0.62cm]{\fm@marca}}}
\fancyhead[R]{\small\bfseries\color{cursfosc}\fm@cursnom}
\renewcommand{\headrule}{\color{curs!45}\hrule height \headrulewidth}
\makeatother

%   \fmcurs{1eso}
% i, al començament del document:
%   \capcalera{Matemàtiques}{1r ESO}{Fitxa d'exercicis · Unitat 2}{Divisibilitat}
% Opcional abans de \input: \def\fmneed{14} (espai mínim abans de cada secció).
% =====================================================================
\makeatletter
\def\fm@logodir{../../logo/pdf}
\providecommand{\fmneed}{10}
\def\fm@cursnom{}

% clau del curs, color i color del text damunt del color
\newcommand{\fm@curs}[3]{\@namedef{fm@color@#1}{#2}\@namedef{fm@text@#1}{#3}}
\fm@curs{1eso}{E6B84A}{navy}
\fm@curs{2eso}{3B82F6}{white}
\fm@curs{3eso}{FF7A50}{white}
\fm@curs{4eso-a}{A78BFA}{white}
\fm@curs{4eso-b}{6D28D9}{white}
\fm@curs{1batx-mi}{9B1C1C}{white}
\fm@curs{1batx-ccss}{EF6B6B}{white}
\fm@curs{2batx-mii}{1E6B47}{white}
\fm@curs{2batx-ccss}{2BB3A3}{white}

\newcommand{\fmcurs}[1]{%
  \@ifundefined{fm@color@#1}{\PackageError{fentmates}{Curs desconegut: #1}{}}{}%
  \edef\fm@tmp{\noexpand\definecolor{curs}{HTML}{\@nameuse{fm@color@#1}}%
               \noexpand\colorlet{curstext}{\@nameuse{fm@text@#1}}}\fm@tmp
  \colorlet{cursfosc}{curs!72!black}%
  \def\fm@logo{\fm@logodir/vertical-#1.pdf}%
  \def\fm@marca{\fm@logodir/marca-#1.pdf}}

% ---------- capçalera de la primera pàgina ----------
% #1 matèria (petit), #2 curs, #3 tipus de document i unitat, #4 títol
\newcommand{\capcalera}[4]{%
  \gdef\fm@cursnom{#2}%
  \par\noindent
  \begin{tcolorbox}[colback=curs,colframe=curs,boxrule=0pt,arc=2.2mm,width=3.5cm,height=2.15cm,
      valign=center,left=3mm,right=2mm,top=1mm,bottom=1mm,nobeforeafter,box align=center]
    \color{curstext}{\footnotesize #1}\par\vspace{1pt}{\fontsize{19}{21}\selectfont\bfseries #2\par}%
  \end{tcolorbox}\hspace{5mm}%
  \begin{minipage}[c]{\dimexpr\linewidth-3.5cm-5mm-2.9cm\relax}
    {\color{cursfosc}\small\bfseries\MakeUppercase{#3}}\par\vspace{2pt}
    {\color{navy}\fontsize{20}{23}\selectfont\bfseries #4\par}
  \end{minipage}\hfill
  \begin{minipage}[c]{2.6cm}\raggedleft\includegraphics[width=2.6cm]{\fm@logo}\end{minipage}\par
  \vspace{2pt}}

% ---------- barres de secció del color del curs ----------
\renewcommand{\seccio}[2]{\par\needspace{\fmneed\baselineskip}\bigskip
  \begin{tcolorbox}[colback=curs,colframe=curs,coltext=curstext,arc=2mm,boxrule=0pt,left=3mm,right=3mm,top=1.2mm,bottom=1.2mm]
  {\large\bfseries #1\quad #2}
  \end{tcolorbox}\vspace{-2pt}}

% ---------- capçalera i peu de les pàgines ----------
% primera pàgina: logos institucionals al peu
\fancypagestyle{primera}{\fancyhf{}\renewcommand{\headrulewidth}{0pt}
  \fancyfoot[L]{\raisebox{-0.32\height}{\includegraphics[height=0.78cm]{logo-gva.png}}\hspace{5mm}%
    \raisebox{-0.32\height}{\includegraphics[height=0.78cm]{logo-vermellar.png}}}
  \fancyfoot[R]{\scriptsize\color{muted}\href{https://creativecommons.org/licenses/by-nc-sa/4.0/deed.ca}{\ccbyncsa}\ \ CC BY-NC-SA 4.0 · fentmates.org\qquad\small\thepage}}
% la resta: el símbol a l'esquerra i el curs a la dreta, amb el títol al centre com abans
\fancyhead[L]{\raisebox{-0.3\height}{\includegraphics[height=0.62cm]{\fm@marca}}}
\fancyhead[R]{\small\bfseries\color{cursfosc}\fm@cursnom}
\renewcommand{\headrule}{\color{curs!45}\hrule height \headrulewidth}
\makeatother

%   \fmcurs{1eso}
% i, al començament del document:
%   \capcalera{Matemàtiques}{1r ESO}{Fitxa d'exercicis · Unitat 2}{Divisibilitat}
% Opcional abans de \input: \def\fmneed{14} (espai mínim abans de cada secció).
% =====================================================================
\makeatletter
\def\fm@logodir{../../logo/pdf}
\providecommand{\fmneed}{10}
\def\fm@cursnom{}

% clau del curs, color i color del text damunt del color
\newcommand{\fm@curs}[3]{\@namedef{fm@color@#1}{#2}\@namedef{fm@text@#1}{#3}}
\fm@curs{1eso}{E6B84A}{navy}
\fm@curs{2eso}{3B82F6}{white}
\fm@curs{3eso}{FF7A50}{white}
\fm@curs{4eso-a}{A78BFA}{white}
\fm@curs{4eso-b}{6D28D9}{white}
\fm@curs{1batx-mi}{9B1C1C}{white}
\fm@curs{1batx-ccss}{EF6B6B}{white}
\fm@curs{2batx-mii}{1E6B47}{white}
\fm@curs{2batx-ccss}{2BB3A3}{white}

\newcommand{\fmcurs}[1]{%
  \@ifundefined{fm@color@#1}{\PackageError{fentmates}{Curs desconegut: #1}{}}{}%
  \edef\fm@tmp{\noexpand\definecolor{curs}{HTML}{\@nameuse{fm@color@#1}}%
               \noexpand\colorlet{curstext}{\@nameuse{fm@text@#1}}}\fm@tmp
  \colorlet{cursfosc}{curs!72!black}%
  \def\fm@logo{\fm@logodir/vertical-#1.pdf}%
  \def\fm@marca{\fm@logodir/marca-#1.pdf}}

% ---------- capçalera de la primera pàgina ----------
% #1 matèria (petit), #2 curs, #3 tipus de document i unitat, #4 títol
\newcommand{\capcalera}[4]{%
  \gdef\fm@cursnom{#2}%
  \par\noindent
  \begin{tcolorbox}[colback=curs,colframe=curs,boxrule=0pt,arc=2.2mm,width=3.5cm,height=2.15cm,
      valign=center,left=3mm,right=2mm,top=1mm,bottom=1mm,nobeforeafter,box align=center]
    \color{curstext}{\footnotesize #1}\par\vspace{1pt}{\fontsize{19}{21}\selectfont\bfseries #2\par}%
  \end{tcolorbox}\hspace{5mm}%
  \begin{minipage}[c]{\dimexpr\linewidth-3.5cm-5mm-2.9cm\relax}
    {\color{cursfosc}\small\bfseries\MakeUppercase{#3}}\par\vspace{2pt}
    {\color{navy}\fontsize{20}{23}\selectfont\bfseries #4\par}
  \end{minipage}\hfill
  \begin{minipage}[c]{2.6cm}\raggedleft\includegraphics[width=2.6cm]{\fm@logo}\end{minipage}\par
  \vspace{2pt}}

% ---------- barres de secció del color del curs ----------
\renewcommand{\seccio}[2]{\par\needspace{\fmneed\baselineskip}\bigskip
  \begin{tcolorbox}[colback=curs,colframe=curs,coltext=curstext,arc=2mm,boxrule=0pt,left=3mm,right=3mm,top=1.2mm,bottom=1.2mm]
  {\large\bfseries #1\quad #2}
  \end{tcolorbox}\vspace{-2pt}}

% ---------- capçalera i peu de les pàgines ----------
% primera pàgina: logos institucionals al peu
\fancypagestyle{primera}{\fancyhf{}\renewcommand{\headrulewidth}{0pt}
  \fancyfoot[L]{\raisebox{-0.32\height}{\includegraphics[height=0.78cm]{logo-gva.png}}\hspace{5mm}%
    \raisebox{-0.32\height}{\includegraphics[height=0.78cm]{logo-vermellar.png}}}
  \fancyfoot[R]{\scriptsize\color{muted}\href{https://creativecommons.org/licenses/by-nc-sa/4.0/deed.ca}{\ccbyncsa}\ \ CC BY-NC-SA 4.0 · fentmates.org\qquad\small\thepage}}
% la resta: el símbol a l'esquerra i el curs a la dreta, amb el títol al centre com abans
\fancyhead[L]{\raisebox{-0.3\height}{\includegraphics[height=0.62cm]{\fm@marca}}}
\fancyhead[R]{\small\bfseries\color{cursfosc}\fm@cursnom}
\renewcommand{\headrule}{\color{curs!45}\hrule height \headrulewidth}
\makeatother

\fmcurs{2batx-ccss}

\begin{document}
\thispagestyle{primera}

\capcalera{Matemàtiques CCSS}{2n Batx}{Solucions de la fitxa · Unitat 4}{Límits i continuïtat}
\vspace{8pt}

\begin{tcolorbox}[colback=soft,colframe=soft,arc=2mm,left=3mm,right=3mm,top=1.5mm,bottom=1.5mm]
\small Intenta resoldre cada exercici abans de mirar-ne la solució. Si no t'ix el mateix resultat, repassa els passos o comprova'l amb la calculadora de límits dels apunts interactius de la unitat.
\end{tcolorbox}

% =====================================================
\seccio{1}{Límits a l'infinit}

\ex{1} \label{ex:g1}\begin{tasks}(2)
  \task $\displaystyle\lim_{x\to+\infty}f(x)=\res{+\infty}$, \ $\displaystyle\lim_{x\to-\infty}f(x)=\res{0}$
  \task $\displaystyle\lim_{x\to+\infty}g(x)=\res{2}$, \ $\displaystyle\lim_{x\to-\infty}g(x)=\res{2}$ \ (asímptota $y=2$)
\end{tasks}

\ex{2} \begin{tasks}(1)
  \task Domini $(0,+\infty)$: \ $\displaystyle\lim_{x\to+\infty}=\res{+\infty}$; \ $\res{\nexists}\displaystyle\lim_{x\to-\infty}$ (la funció no està definida per a $x\leq0$).
  \task Domini $[-3,+\infty)$: \ $\displaystyle\lim_{x\to+\infty}=\res{+\infty}$; \ $\res{\nexists}\displaystyle\lim_{x\to-\infty}$.
  \task Domini $[-2,2]$: no arriba a cap infinit, $\res{\nexists}$ cap dels dos límits.
\end{tasks}

\ex{1} \begin{tasks}(4)
  \task $\lim 3x^4=\res{+\infty}$
  \task $\lim(-x^3)=\res{+\infty}$
  \task $\lim 2x^5=\res{-\infty}$
  \task $\lim(-2x^2)=\res{-\infty}$
\end{tasks}

\ex{1} \begin{tasks}(3)
  \task Mateix grau: $\res{\frac25}$
  \task Grau de baix major: $\res{0}$
  \task $\lim\frac{4x^2}{3}=\res{+\infty}$
  \task $\lim(-3x)=\res{+\infty}$
  \task $\frac{7}{-1}=\res{-7}$
  \task $\lim\frac{x}{100}=\res{+\infty}$
\end{tasks}

\ex{2} \begin{tasks}(3)
  \task $\lim\frac{3x}{2x}=\res{\frac32}$
  \task $\lim\frac{x^{3/2}}{x}=\lim\sqrt x=\res{+\infty}$
  \task $\lim\frac{2x}{x}=\res{2}$
  \task $3^{-\infty}+1=\res{1}$
  \task $0{,}4^{-\infty}=\res{+\infty}$
  \task $\left(\tfrac52\right)^{-\infty}=\res{0}$
\end{tasks}

\ex{3} $\sqrt{4x^2+3}\approx\sqrt{4x^2}=2|x|$. \ Quan $x\to+\infty$, $|x|=x$: \ $\lim\frac{6x}{2x}=\res{3}$. \ Quan $x\to-\infty$, $|x|=-x$: \ $\lim\frac{6x}{-2x}=\res{-3}$. L'arrel sempre és positiva, mentre que el numerador canvia de signe.

\ex{3} \begin{tasks}(2)
  \task $\lim\frac{x^2}{2x^2}=\res{\frac12}$
  \task $\sqrt{x^2}=|x|=-x$: \ $\lim\frac{-x}{3x}=\res{-\frac13}$
  \task Només està definida per a $x\geq-1$: \ $\res{\nexists}$
  \task Si $x\to-\infty$, $x^3-5x+3<0$: l’arrel no existeix, $\res{\nexists}$
\end{tasks}


% =====================================================
\seccio{2}{Comparació d'infinits}

\ex{1} $\res{\ln x\ll\sqrt x\ll x^3\ll x^{10}\ll 2^x\ll e^x}$

\ex{2} \begin{tasks}(3)
  \task $\res{0}$
  \task $\lim\frac{2^x}{3\cdot2^x}=\res{\frac13}$
  \task $\res{0}$
  \task $\res{-\infty}$
  \task $\res{+\infty}$
  \task $\lim\frac{3^x}{3^x}=\res{1}$
\end{tasks}


% =====================================================
\seccio{3}{Indeterminacions}

\ex{2} \begin{tasks}(1)
  \task $\dfrac{x^2-x(x+1)}{x+1}=\dfrac{-x}{x+1}\to\res{-1}$
  \task $\dfrac{2x^2+1-2x(x-1)}{x-1}=\dfrac{2x+1}{x-1}\to\res{2}$
  \task $\dfrac{x^3(x+1)-x^2(x^2+1)}{(x^2+1)(x+1)}=\dfrac{x^3-x^2}{x^3+x^2+x+1}\to\res{1}$
\end{tasks}

\ex{2} \begin{tasks}(1)
  \task $\dfrac{8x}{\sqrt{x^2+8x}+x}\to\dfrac{8x}{2x}=\res{4}$
  \task $\dfrac{x}{\sqrt{4x^2+x}+2x}\to\dfrac{x}{4x}=\res{\frac14}$
  \task $\dfrac{4x}{\sqrt{x^2+3x}+\sqrt{x^2-x}}\to\dfrac{4x}{2x}=\res{2}$
  \task $\sqrt{x^2+1}\approx x$, i $x-3x=-2x\to\res{-\infty}$ (no hi ha indeterminació real).
\end{tasks}

\ex{2} \begin{tasks}(2)
  \task $e^{\lim\frac5x\cdot x}=\res{e^5}$
  \task $e^{\lim\frac{3}{x-1}\cdot2x}=\res{e^6}$
  \task $e^{\lim\frac{-4}{2x+3}(x+1)}=\res{e^{-2}}$
  \task $e^{\lim\frac{4}{x^2-1}\cdot x^2}=\res{e^4}$
  \task $3^{+\infty}=\res{+\infty}$ (no és $1^\infty$)
  \task $e^{\lim\frac{-4}{x}\cdot3x}=\res{e^{-12}}$
\end{tasks}

\ex{2} \begin{tasks}(2)
  \task $\lim\frac{x+3x}{2x}=\res{2}$
  \task $\lim\frac{x^2}{x^2}=\res{1}$
\end{tasks}

\ex{2} \pau{Aragó, EvAU extraordinària 2020}Per a $x>4$: $+\infty-\infty$. Pel conjugat: \ $\dfrac{4x^2-7x-4x^2}{\sqrt{4x^2-7x}+2x}=\dfrac{-7x}{\sqrt{4x^2-7x}+2x}\to\dfrac{-7x}{4x}=\res{-\frac74}$.

\ex{2} \pau{Balears, PBAU extraordinària 2021}\begin{tasks}(1)
  \task $P(0)=\res{2}$ milions; \ $P(5)=\frac{32}{36}=\frac89\approx\res{0{,}889}$ milions ($\approx888\,889$ individus).
  \task $\frac{2+t+t^2}{(t+1)^2}<1\iff2+t+t^2<t^2+2t+1\iff t>1$: \ \rest{a partir del primer any}.
  \task $\displaystyle\lim_{t\to+\infty}P(t)=\res{1}$: tendeix a \rest{un milió} d'individus.
\end{tasks}

\ex{2} \pau{Castella i Lleó, EBAU juny 2017}$\displaystyle\lim_{x\to+\infty}\left(\frac{90x-45}{x+8}+27\right)=90+27=117$: \ el benefici s'acosta a \rest{117\,000\,€}.

\ex{3} \pau{Catalunya, PAU setembre 2018}Asímptota vertical $x=1$: el denominador s'anul·la en 1, $c+1=0\Rightarrow c=-1$. \ Horitzontal: $\displaystyle\lim_{x\to\pm\infty}\frac{ax+b}{cx+1}=\frac ac=2\Rightarrow a=-2$. \ Punt: $\frac{-4+b}{-2+1}=-5\Rightarrow b=9$. \ $\res{a=-2,\ b=9,\ c=-1}$: \ $f(x)=\frac{-2x+9}{1-x}$.

\ex{2} \begin{tasks}(1)
  \task $(-\infty)-(+\infty)=\res{-\infty}$: no hi ha indeterminació.
  \task $(-\infty)-(-\infty)$, IND. \ $\dfrac{3x^2(x^2+1)-x^3(x+1)}{(x+1)(x^2+1)}=\dfrac{2x^4-x^3+3x^2}{x^3+x^2+x+1}\to\lim 2x=\res{-\infty}$
  \task Conjugat: $\dfrac{x-1}{\sqrt{x^2+x}+\sqrt{x^2+1}}\to\dfrac{x}{2|x|}=\dfrac{x}{-2x}=\res{-\frac12}$
  \task $(-\infty)-(+\infty)=\res{-\infty}$: no hi ha indeterminació.
\end{tasks}

\ex{2} \begin{tasks}(1)
  \task La base tendeix a $\frac32\neq1$: \ $\left(\frac32\right)^{-\infty}=\res{0}$
  \task $1^{-\infty}$: \ $\frac{2x+1}{2x-1}-1=\frac{2}{2x-1}$; \ $\lim\frac{2(x-1)}{2x-1}=1\Rightarrow\res{e}$
  \task $1^{-\infty}$: \ $\frac{x^2+1}{x^2-1}-1=\frac{2}{x^2-1}$; \ $\lim\frac{2(x^2-4)}{(x^2-1)(x-1)}=0\Rightarrow e^0=\res{1}$
\end{tasks}


% =====================================================
\seccio{4}{Límits en un punt i indeterminació 0/0}

\ex{1} \begin{tasks}(2)
  \task $\lim_{x\to-1^-}f=+\infty$, \ $\lim_{x\to-1^+}f=-\infty$: \ $\res{\nexists}$
  \task $\lim_{x\to1^-}f=\lim_{x\to1^+}f=+\infty$: \ $\res{+\infty}$
  \task $\res{0}$
  \task $\lim_{x\to-2^-}g=+\infty$, \ $\lim_{x\to-2^+}g=-\infty$: \ $\res{\nexists}$
  \task $\lim_{x\to2^-}g=-\infty$, \ $\lim_{x\to2^+}g=+\infty$: \ $\res{\nexists}$
  \task $\res{0}$
\end{tasks}

\ex{1} \label{ex:g4}$\displaystyle\lim_{x\to+\infty}f(x)=\res{1}$, \ $\displaystyle\lim_{x\to-\infty}f(x)=\res{+\infty}$. \ $\lim_{x\to0^-}f=-2$ i $\lim_{x\to0^+}f=2$: \ $\res{\nexists}\displaystyle\lim_{x\to0}f(x)$.

\ex{1} \begin{tasks}(3)
  \task $8-4+1=\res{5}$
  \task $\frac{4}{-2}=\res{-2}$
  \task $\frac60$: \ $\res{-\infty}$ per l’esquerra, $\res{+\infty}$ per la dreta: $\nexists$
  \task $\frac{-1}{0^+}$: \ $\res{-\infty}$ pels dos costats
  \task $\sqrt{0^+}=\res{0}$
  \task $\frac20$: \ $\res{-\infty}$ per l’esquerra, $\res{+\infty}$ per la dreta: $\nexists$
\end{tasks}

\ex{2} $x=2$: $\lim_{x\to2^-}2x=4$, $\lim_{x\to2^+}\frac8x=4\Rightarrow\res{\lim_{x\to2}f(x)=4}$. \ \ $x=-1$: $\lim_{x\to-1^-}(x^2-3)=-2$, $\lim_{x\to-1^+}2x=-2\Rightarrow\res{\lim_{x\to-1}f(x)=-2}$.

\ex{2} \begin{tasks}(2)
  \task $\frac{(x-2)(x-3)}{(x+3)(x-3)}=\frac{x-2}{x+3}\to\res{\frac16}$
  \task $\frac{(x-2)(x^2+2x+4)}{x-2}\to\res{12}$
  \task $\frac{(x-1)^2(x+1)}{(x-1)(x+1)}=x-1\to\res{0}$
  \task $\frac{(x-1)(x+1)(x^2+1)}{x(x-1)}\to\frac{2\cdot2}{1}=\res{4}$
  \task $\frac{x(x+3)}{(x+3)^2}=\frac{x}{x+3}\to\frac{-3}{0}$: \ $\res{+\infty}$ per l’esquerra i $\res{-\infty}$ per la dreta, $\nexists$
  \task $\frac{(x+5)-9}{(x-4)(\sqrt{x+5}+3)}=\frac{1}{\sqrt{x+5}+3}\to\res{\frac16}$
\end{tasks}

\ex{2} \begin{tasks}(2)
  \task $\frac00$: \ $\frac{x-3}{(x-3)^2(x+1)}=\frac{1}{(x-3)(x+1)}\to\frac10$: \ $\res{-\infty}$ per l’esquerra, $\res{+\infty}$ per la dreta, $\nexists$
  \task $\frac{2x-(x+1)}{x(x-2)}=\frac{x-1}{x(x-2)}\to\frac10$: \ $\res{-\infty}$ per l’esquerra, $\res{+\infty}$ per la dreta, $\nexists$
  \task $\frac{0}{1}=\res{0}$ \ (no hi ha indeterminació)
  \task $1^{0}=\res{1}$ \ (no hi ha indeterminació)
  \task Ruffini: $\frac{(x+1)(x^2-3x+5)}{(x+1)(x-7)}\to\frac{9}{-8}=\res{-\frac98}$
  \task $\frac{(x-2)(2x+5)}{(x-2)^2(x+1)}=\frac{2x+5}{(x-2)(x+1)}\to\frac90$: \ $\res{-\infty}$ per l’esquerra, $\res{+\infty}$ per la dreta, $\nexists$
\end{tasks}


% =====================================================
\seccio{5}{Continuïtat en un punt. Funcions racionals}

\ex{1} $\lim_{x\to1^-}f=\lim_{x\to1^+}f=+\infty$: el límit és $+\infty$ (no és un nombre) i $\nexists f(1)$: \rest{discontinuïtat asimptòtica}. \ A l'exercici \ref{ex:g4}: $f(0)=2$, laterals $-2$ i $2$: \rest{salt finit} en $x=0$.

\ex{1} \begin{tasks}(2)
  \task $\nexists f(-2)$, $\lim(x-2)=-4$: \ \rest{evitable}
  \task Laterals $-\infty$ i $+\infty$: \ \rest{asimptòtica}
  \task $h(0)=1$, laterals $1$ i $1$: \ \rest{contínua}
  \task $k(1)=2$, laterals $2$ i $4$: \ \rest{salt finit}
\end{tasks}

\ex{2} $x^2-4x+3=0\Rightarrow x=1,\ x=3$. \ $f(x)=\frac{(x-1)(x+2)}{(x-1)(x-3)}$. \ $x=1$: $\lim\frac{x+2}{x-3}=-\frac32$ \ $\Rightarrow$ \rest{evitable}. \ $x=3$: $\frac50$ \ $\Rightarrow$ \rest{asimptòtica}. Contínua en $\mathbb R-\{1,3\}$.

\ex{2} \pau{Cantàbria, EBAU juliol 2019}$x^2-2x-15=(x-5)(x+3)$ i $x^2-x-12=(x-4)(x+3)$: discontínua en $\res{x=-3}$ i $\res{x=5}$. \ En $x=-3$: $\lim\frac{x-4}{x-5}=\frac{-7}{-8}=\frac78$ \ (evitable): definint $\res{f(-3)=\frac78}$ deixa de ser discontínua. En $x=5$ ($\frac{1}{0}$) la discontinuïtat és asimptòtica i no es pot evitar.

\ex{2} \pau{Cantàbria, EBAU ordinària 2020}\begin{tasks}(1)
  \task $x^2+5x+6=(x+2)(x+3)$: discontínua en $\res{x=-2}$ i $\res{x=-3}$.
  \task $f(x)=\frac{2(x+2)}{(x+2)(x+3)}$; \ $\lim_{x\to-2}\frac{2}{x+3}=2$: \ definint $\res{f(-2)=2}$ s'evita la discontinuïtat de $x=-2$.
  \task $\lim_{x\to-3^-}f(x)=\res{-\infty}$, \ $\lim_{x\to-3^+}f(x)=\res{+\infty}$: \ $x=-3$ és una asímptota vertical; la gràfica baixa per l'esquerra i puja per la dreta.
\end{tasks}

\ex{2} \pau{Navarra, EvAU ordinària 2021}És racional: contínua excepte on $2x-2=0$, és a dir, en $x=1$. \ $\lim_{x\to1}f(x)=\frac{8}{0}$: \ $\lim_{x\to1^-}=-\infty$, \ $\lim_{x\to1^+}=+\infty$. \ Contínua en $\mathbb R-\{1\}$; en $x=1$, \rest{discontinuïtat asimptòtica} (inevitable, de salt infinit segons la terminologia de Navarra).


% =====================================================
\seccio{6}{Continuïtat de funcions a trossos}

\ex{1} \pau{Navarra, EvAU extraordinària 2021}Trossos polinòmics: continus. \ En $x=2$: $f(2)=0$, \ $\lim_{x\to2^-}(x^2-2x)=0$, \ $\lim_{x\to2^+}(-x^2+6x-8)=0$. \ \rest{Contínua en tot $\mathbb R$}.

\ex{2} \pau{Castella i Lleó, EBAU setembre 2017}El primer tros no està definit en $x=\pm2$, que no són de $(0,2)$: continu. \ En $x=2$: $f(2)=4$, \ $\lim_{x\to2^-}\frac{1}{4-x^2}=\frac{1}{0^+}=+\infty$, \ $\lim_{x\to2^+}2x=4$. \ Contínua en $\res{(0,2)\cup(2,4)}$; en $x=2$, \rest{salt infinit}.

\ex{2} \pau{Madrid, EvAU juny 2018}$f(2)=\frac41=4$, \ $\lim_{x\to2^-}f=4$, \ $\lim_{x\to2^+}f=\frac{12-4}{4}=2$: \ \rest{no és contínua en $x=2$} (salt finit). \ A més, el primer tros no existeix en $x=1$ ($\frac30$): \rest{asimptòtica en $x=1$}. ($x=-2$ no és de $x>2$.) \ Contínua en $\mathbb R-\{1,2\}$.

\ex{2} \pau{Madrid, EvAU juliol 2018}El denominador $3+x$ s'anul·la en $-3$, que no és de $x\geq0$: domini $\res{\mathbb R}$. \ En $x=0$: $f(0)=\frac23$, \ $\lim_{x\to0^-}(x^3+2e^x)=2$, \ $\lim_{x\to0^+}\frac{2}{3+x}=\frac23$: \ salt finit. \ Contínua en $\res{\mathbb R-\{0\}}$.

\ex{2} \pau{Aragó, EvAU extraordinària 2020}Trossos: el primer, continu en $x<-1$; el segon, polinòmic; el tercer, $4x^2-7x=x(4x-7)>0$ si $x>4$: continu. \ $x=-1$: $\lim_{x\to-1^-}\frac{3}{x+1}=-\infty$, $\lim_{x\to-1^+}=-1-4-2-10=-17=f(-1)$: \ salt infinit. \ $x=4$: $f(4)=64-64+8-10=-2$, $\lim_{x\to4^+}=\sqrt{36}-8=-2$: contínua. \ \rest{Contínua en $\mathbb R-\{-1\}$}.

\ex{3} \pau{Madrid, EvAU juny 2019}$x=0$: $f(0)=\lim_{x\to0^-}f=1+k$, \ $\lim_{x\to0^+}(1-x^2)=1$: contínua només si $k=0$. \ $x=3$: $f(3)=-8$, \ $\lim_{x\to3^+}\frac{1}{x-3}=+\infty$: sempre discontínua (salt infinit). \ \rest{Si $k=0$, contínua en $\mathbb R-\{3\}$; si $k\neq0$, contínua en $\mathbb R-\{0,3\}$}.

\ex{2} \begin{tasks}(1)
  \task $f$: \ $x=0$: $\nexists f(0)$ i els laterals valen $1$ i $1$: \rest{evitable}. \ $x=1$: $f(1)=-1$, $\lim_{x\to1^-}f=2$: \rest{salt finit}. Contínua en $\mathbb R-\{0,1\}$.
  \task $g$: \ $x=3$: $g(3)=0$, $\lim_{x\to3^+}g=0$: contínua. \ $x=6$: $\lim_{x\to6^-}g=3$, $g(6)=0$: \rest{salt finit}. Contínua en $\mathbb R-\{6\}$.
\end{tasks}


% =====================================================
\seccio{7}{Funcions a trossos amb paràmetres}

\ex{1} \pau{Madrid, EvAU setembre 2017}$\lim_{x\to-1^-}(ax+1)=-a+1$, \ $\lim_{x\to-1^+}=f(-1)=1-1-2=-2$. \ $-a+1=-2\Rightarrow\res{a=3}$.

\ex{1} \pau{Andalusia, PEvAU juny 2019}El primer tros només falla en $x=1$, que no és de $x<0$. \ $\lim_{x\to0^-}\frac{1}{x-1}=-1$, \ $\lim_{x\to0^+}(x^2+a)=a=f(0)$ \ $\Rightarrow\res{a=-1}$.

\ex{2} \pau{Castella-la Manxa, EvAU extraordinària 2020}$f(0)=\lim_{x\to0^-}f=-t$, \ $\lim_{x\to0^+}f=t^2-5t+4$. \ $t^2-5t+4=-t\Rightarrow t^2-4t+4=0\Rightarrow(t-2)^2=0\Rightarrow\res{t=2}$.

\ex{2} \pau{Castella-la Manxa, EvAU extraordinària 2020}$f(c)=\lim_{x\to c^-}f=-4c+4$, \ $\lim_{x\to c^+}f=c^2-4c$. \ $c^2-4c=-4c+4\Rightarrow c^2=4\Rightarrow\res{c=-2\ \text{o}\ c=2}$.

\ex{2} \pau{Balears, PBAU extraordinària 2021}$x=-2$: $8+a=4-5=-1\Rightarrow\res{a=-9}$. \ $x=1$: $\lim_{x\to1^-}(x^2-5)=-4$ i $f(1)=\lim_{x\to1^+}(-7x+3)=-4$: contínua per a qualsevol $a$.

\ex{2} \pau{Catalunya, PAU setembre 2018}$x=-1$: $1=-a+b$. \ $x=2$: $2a+b=4$. \ Restant: $3a=3\Rightarrow\res{a=1,\ b=2}$.

\ex{3} \pau{Aragó, EvAU setembre 2018}$x=0$: $\lim_{x\to0^-}\frac{x+b}{x^2+1}=b$, \ $f(0)=4\Rightarrow b=4$. \ $x=-2$: $-2a+1=\frac{-2+4}{5}=\frac25\Rightarrow\res{a=\frac{3}{10},\ b=4}$.

\ex{3} \pau{Cantàbria, EBAU ordinària 2020}$x=-1$: $a-2-1=1-5\Rightarrow a-3=-4\Rightarrow a=-1$. \ $x=3$: $f(3)=4=\frac{b+3}{7}\Rightarrow b=25$. \ $\res{a=-1,\ b=25}$ \ ($3x-2=0$ en $x=\frac23$, fora de $x>3$).

\ex{3} \pau{Cantàbria, EBAU juliol 2019}$x=-2$: $15=4-2a+5\Rightarrow a=-3$. \ $x=0$: $5=\frac{15}{b}\Rightarrow b=3$. \ $\res{a=-3,\ b=3}$ \ (amb $b=3$ el denominador $x+3$ no s'anul·la en $[0,3)$).

\ex{2} \begin{tasks}(1)
  \task $\lim_{x\to1^-}\left(-\frac{1}{x^2}\right)=-1$, \ $f(1)=1+a$ \ $\Rightarrow\res{a=-2}$.
  \task El primer tros no està definit en $x=0$: $\lim_{x\to0}\left(-\frac{1}{x^2}\right)=-\infty$ pels dos costats, \rest{asimptòtica}. Contínua en $\mathbb R-\{0\}$.
\end{tasks}

\ex{3} \begin{tasks}(1)
  \task $f$: \ $x=-1$: $1+a=7\Rightarrow a=6$; \ $x=1$: $7=-1+b\Rightarrow b=8$. \ $\res{a=6,\ b=8}$ ($x=0$ no és de $x\leq-1$).
  \task $g$: \ $x=-2$: $4+a=-1\Rightarrow a=-5$; \ $x=1$: $-4=b+3\Rightarrow b=-7$. \ $\res{a=-5,\ b=-7}$.
\end{tasks}


% =====================================================
\seccio{8}{Asímptotes}

\ex{1} \begin{tasks}(1)
  \task AV $\res{x=1}$ ($\lim_{x\to1^-}f=-\infty$, $\lim_{x\to1^+}f=+\infty$); \ AH $\res{y=3}$.
  \task AV $\res{x=-4}$ ($+\infty$ per l’esquerra, $-\infty$ per la dreta) i $\res{x=4}$ ($-\infty$ per l’esquerra, $+\infty$ per la dreta); \ AH $\res{y=0}$.
  \task Cap AV ($x^2+3\neq0$); \ AH $\res{y=2}$.
  \task AV $\res{x=2}$ ($\frac{3}{0^+}$: $+\infty$ pels dos costats); \ AH $\res{y=0}$.
  \task $x^2+x-2=(x-1)(x+2)$. En $x=1$: $\frac00$, $\frac{-(x-1)(x+1)}{(x-1)(x+2)}=\frac{-(x+1)}{x+2}\to-\frac23$: no hi ha AV. \ AV $\res{x=-2}$ ($\frac{-3}{0}$: $-\infty$ per l’esquerra, $+\infty$ per la dreta); \ AH $\res{y=-1}$.
  \task $\frac{(x+3)(x-3)}{x+3}=x-3$: en $x=-3$ hi ha un forat, \rest{no té AV}. Graus $2>1$: \rest{no té AH} (la gràfica és la recta $y=x-3$ sense el punt $(-3,-6)$).
\end{tasks}

\ex{2} \begin{tasks}(1)
  \task $\lim_{x\to-\infty}(3^x+2)=0+2$: \ AH $\res{y=2}$ per l’esquerra; per la dreta $+\infty$, no n’hi ha.
  \task $\lim_{x\to+\infty}(4-e^{-x})=4-0$: \ AH $\res{y=4}$ per la dreta; per l’esquerra $-\infty$, no n’hi ha.
  \task $x\to+\infty$: $\lim\frac{6\cdot2^x}{2^x}=6$, \ AH $\res{y=6}$. \ $x\to-\infty$: $2^x\to0$, $\frac{0}{0+1}=0$, \ AH $\res{y=0}$. Té una AH diferent per cada costat.
\end{tasks}

\ex{2} \begin{tasks}(1)
  \task $x^2+2=(x-1)(x+1)+3$: \ AO $\res{y=x+1}$, \ AV $\res{x=1}$. \ $\frac{3}{x-1}>0$ si $x\to+\infty$ (per damunt) i $<0$ si $x\to-\infty$ (per davall).
  \task $3x^2-x=(x+2)(3x-7)+14$: \ AO $\res{y=3x-7}$, \ AV $\res{x=-2}$. \ $\frac{14}{x+2}$: per damunt quan $x\to+\infty$ i per davall quan $x\to-\infty$.
  \task $f(x)=x+\frac{2}{x^2}$: \ AO $\res{y=x}$, \ AV $\res{x=0}$ ($+\infty$ pels dos costats). \ $\frac{2}{x^2}>0$: sempre per damunt.
  \task $-x^2+4x=(x-3)(-x+1)+3$: \ AO $\res{y=-x+1}$, \ AV $\res{x=3}$. \ $\frac{3}{x-3}$: per damunt quan $x\to+\infty$ i per davall quan $x\to-\infty$.
\end{tasks}

\ex{2} \begin{tasks}(1)
  \task $\frac{x(x-1)}{(x-1)(x+1)}$: en $x=1$, $\frac x{x+1}\to\frac12$ (forat, no és AV). \ AV $\res{x=-1}$; \ AH $\res{y=1}$.
  \task AV $\res{x=-2}$ i $\res{x=2}$; \ $x^3=(x^2-4)\cdot x+4x$: \ AO $\res{y=x}$.
  \task AV $\res{x=1}$. Graus 3 i 1: el quocient $x^2+x+1$ és de grau 2, \rest{no té AH ni AO}.
\end{tasks}

\ex{3} \begin{tasks}(1)
  \task AV en $x=-2$: el denominador s'anul·la en $-2$ $\Rightarrow\res{b=2}$ (el numerador $4a+1\neq0$). \ El quocient de $(ax^2+1):(x+2)$ comença per $ax$: pendent $\res{a=3}$.
  \task $3x^2+1=(x+2)(3x-6)+13$: \ AO $\res{y=3x-6}$.
\end{tasks}

\ex{2} \begin{tasks}(1)
  \task $V(6)=\frac{300}{10}=\res{30}$ milers d’unitats.
  \task $\displaystyle\lim_{t\to+\infty}\frac{50t}{t+4}=\res{50}$: AH $y=50$. A llarg termini les vendes s’acosten a \rest{50\,000 unitats}, sense arribar-hi.
\end{tasks}


\end{document}
